\PassOptionsToPackage{dvipsnames}{xcolor}
%\documentclass[12pt]{beamer}
%\documentclass[10pt]{beamer}
\documentclass[10pt, handout]{beamer}

\usepackage{mathtools} %% previously \usepackage{amsmath}
\usepackage{amsfonts}
\usepackage{amssymb} % http://www.ctan.org/tex-archive/macros/latex/required/psnfss/psnfss2e.pdf
%\usepackage[official]{eurosym}
\usepackage{helvet}
\usepackage{natbib}
\usepackage{lscape, longtable, array, booktabs, multirow, tabularx, graphics, epstopdf}
\usepackage{setspace}
\usepackage{listings}
\usepackage {tikz}

\usefonttheme[onlymath]{serif}

\usepackage{accents}
\newcommand*{\dt}[1]{%
   \accentset{\mbox{\LARGE $.$}}{#1}}

%\usepackage[T1]{fontenc}

\usetheme{Boadilla}
%\usecolortheme{dolphin}
%\useoutertheme{infolines}
\definecolor{dollarbill}{rgb}{0.52, 0.73, 0.4}
%\setbeamercovered{transparent}% Dim out "inactive" elements
\setbeamercovered{invisible}% Default
\beamertemplatenavigationsymbolsempty
\setbeamerfont{page number in head/foot}{size=\scriptsize}
\setbeamertemplate{footline}[frame number]


\usepackage{hyperref} 
\hypersetup{colorlinks=true}
%\hypersetup{linkcolor=RoyalBlue, urlcolor=ForestGreen, citecolor=Magenta}
\hypersetup{allcolors=RoyalBlue} %% 

%\renewcommand{\thefootnote}{\fnsymbol{footnote}}
\renewcommand{\thefootnote}{\arabic{footnote}}

\DeclareMathOperator*{\argmax}{argmax} % thin space, limits underneath in displays


\title[Eigenvalues \& Eigenvectors]{Eigenvalues \& Eigenvectors}
\vskip 0.1in

\subtitle{Mathematical Methods for Economics (771)}
%\author[H Hollander]{Hylton Hollander}

\institute[Stellenbosch University]{\large{
Stellenbosch University}}


%\date[April, 2020]{April, 2020}

\date{}

\begin{document}


\frame{\titlepage}

\frame{\frametitle{Contents:} \tableofcontents}


\section{Readings}

\frame{
    \frametitle{Readings} \vskip 0.1in


\begin{itemize}
\item Simon \& Blume, Chapter 23.1-23.3, 23.6-23.7 (recommended self-study 23.8-23.9) \vskip0.1in

\item Additional: \href{https://www.youtube.com/playlist?list=PLZHQObOWTQDPD3MizzM2xVFitgF8hE_ab}{3Blue1Brown}

\end{itemize}

}



\section{Review of Linear Independence (Chapter 11)}


\subsection{Linear combinations and span}

\begin{frame}
\frametitle{Linear combinations and span} \vskip 0.1in

The set of all scalar multiplies of a nonzero vector $\mathbf{v}$ is a straight line through the origin. \\ \vskip0.1in

Formally, we denote this set by:
\begin{equation}
\mathcal{L}[\mathbf{v}] \equiv \{ r. \mathbf{v} : r \in \mathbf{R} \} ~,
\end{equation}
and call it the line \textit{generated} or \textit{spanned} by  $\mathbf{v}$.

\begin{figure}
\centering
\includegraphics[width=0.6\textwidth]{line-span.png}
\caption{The line $\mathcal{L}$ spanned by vector $\mathbf{v}$}
\end{figure}


\end{frame}


\begin{frame}
\frametitle{Linear combinations and span} \vskip 0.1in

The set spanned by two nonzero vectors $\mathbf{v_1}$ and $\mathbf{v_2}$ is given by:

\begin{equation}
\mathcal{L}[\mathbf{v_1},\mathbf{v_2}] \equiv  \{ r_1. \mathbf{v_1} + r_2. \mathbf{v_2} : r_1 , r_2 \in \mathbf{R} \} ~,
\end{equation}
If $\mathbf{v_1}$ is a multiple of $\mathbf{v_2}$, or vice versa, we say $\mathbf{v_1}$ and $\mathbf{v_2}$ are \textbf{linearly dependent}. Otherwise, we say that $\mathbf{v_1}$ and $\mathbf{v_2}$ are \textbf{linearly \textit{in}dependent}.

\begin{figure}
\centering
\includegraphics[width=0.6\textwidth]{line-multi.png}
\caption{If $\mathbf{v_1}$ is a multiple of $\mathbf{v_2}$, then $\mathcal{L}[\mathbf{v_1},\mathbf{v_2}] = \mathcal{L}[\mathbf{v_2}]$ is simply a line spanned by $\mathbf{v_2}$.}
\end{figure}

 
\end{frame}



\begin{frame}
\frametitle{Linear combinations and span} \vskip 0.1in

%Otherwise, we say that $\mathbf{v_1}$ and $\mathbf{v_2}$ are \textbf{linearly \textit{in}dependent}.  \\ \vskip0.1in

That is, given the linear combination of two nonzero vectors $\mathbf{v_1}$ and $\mathbf{v_2}$, $c_{1}\mathbf{v_1} + c_{2}\mathbf{v_2} = 0$, we say $\mathbf{v_1}$ and $\mathbf{v_2}$ are:
\begin{itemize}
\item \textbf{linearly dependent} if $c_1$ or $c_2$ nonzero ($\mathcal{L}[\cdot]$ is a line);
\item \textbf{linearly \textit{in}dependent} if $c_1 = c_2 = 0$ ($\mathcal{L}[\cdot]$ is a plane).
\end{itemize}
\begin{definition}
Vectors $\mathbf{v_1}, \mathbf{v_2}, \ldots , \mathbf{v_k}$ in $\mathbf{R^n}$ are \textbf{linearly dependent} if and only if there exists scalars $c_1, c_2, \ldots, c_k$,\textit{ not all zero}, such that $$ c_1\mathbf{v_1} + c_2\mathbf{v_2} + \ldots + c_k\mathbf{v_k} = 0$$
Vectors $\mathbf{v_1}, \mathbf{v_2}, \ldots , \mathbf{v_k}$ in $\mathbf{R^n}$ are \textbf{linearly independent} if and only if $c_1\mathbf{v_1} + c_2\mathbf{v_2} + \ldots + c_k\mathbf{v_k} = 0$ for scalars $c_1, \ldots, c_k$ implies that $c_1 = \cdots = c_k = 0$.
\end{definition}

\end{frame}



\begin{frame}
%\frametitle{Linear combinations and span} \vskip 0.1in

\begin{theorem}[11.1]
Vectors $\mathbf{v_1}, \mathbf{v_2}, \ldots , \mathbf{v_{\color{red}k}}$ in $\mathbf{R^{\color{Green}n}}$ are \textbf{linearly dependent} if and only if the linear system $$ A \begin{pmatrix} c_1 \\ \vdots \\ c_{\color{red}k} \end{pmatrix} = 0 $$
has a nonzero solution ($c_1, c_2, \ldots, c_{\color{red}k}$), where A is the ${\color{Green}n} \times {\color{red}k}$ matrix whose columns are the vectors $\mathbf{v_1}, \ldots , \mathbf{v_{\color{red}k}}$ under study: 
$$A = \begin{pmatrix} \mathbf{v_1} & \mathbf{v_2} & \cdots & \mathbf{v_{\color{red}k}} \end{pmatrix} ~.$$
\end{theorem}

\begin{theorem}[11.2, case $k = n$]
A set of $n$ vectors $\mathbf{v_1},  \ldots , \mathbf{v_n}$ in $\mathbf{R^n}$ are \textbf{linearly independent} if and only if $$\det{A} = \det \begin{pmatrix} \mathbf{v_1} & \mathbf{v_2} & \cdots & \mathbf{v_n} \end{pmatrix} \neq 0 ~.$$
\end{theorem}
%% uses the fact that a square matrix is nonsingular iff its determinant is not zero 
\begin{theorem}[11.3, case $k > n$]
If $k > n$, any set A set of $k$ vectors in $\mathbf{R^n}$ is \textbf{linearly dependent}.
\end{theorem}

\end{frame}






\subsection{Spanning sets}

\begin{frame}
\frametitle{Spanning sets} \vskip 0.1in

Recall that the set spanned by  $\mathbf{v_1}, \mathbf{v_2}, \ldots , \mathbf{v_{\color{red}k}}$ can be written as $$\mathcal{L}[ \mathbf{v_1}, \mathbf{v_2}, \ldots , \mathbf{v_{k}}] \equiv \{ c_1\mathbf{v_1} + c_2\mathbf{v_2} + \ldots + c_k\mathbf{v_k} : c_1, c_2, \ldots, c_k \in \mathbf{R} \} ~.$$ \\ \vskip0.1in

Suppose we are given a subset $V$ of $\mathbf{R^n}$. Then there exists
$\mathbf{v_1}, \mathbf{v_2}, \ldots , \mathbf{v_{k}}$ in $\mathbf{R^n}$ such that every vector in $V$ can be written as a linear combination of $\mathbf{v_1}, \mathbf{v_2}, \ldots , \mathbf{v_{k}}$: $$V = \mathcal{L}[\mathbf{v_1}, \mathbf{v_2}, \ldots , \mathbf{v_{k}}] ~.$$
That is, $\mathbf{v_1}, \mathbf{v_2}, \ldots , \mathbf{v_{k}}$ \textbf{span} $V$. 


\end{frame}

\begin{frame}
\frametitle{Spanning sets} \vskip 0.1in



\begin{example}[11.4]
The $x_{1}x_{2}$-plane in $\mathbf{R^3}$ is the span of the unit vectors $e_1 = (1, 0, 0)$ and $e_2 = (0, 1, 0)$, because any vector $(a,b,0)$ in this plane can be written as a linear combination of $e_1$ and $e_2$:
$$ \begin{pmatrix} a \\ b \\ 0 \end{pmatrix} = a \begin{pmatrix} 1 \\ 0 \\ 0 \end{pmatrix} + b \begin{pmatrix} 0 \\ 1 \\ 0 \end{pmatrix} $$
\end{example}

\begin{theorem}[11.6]
A set of $k$ vectors that spans $\mathbf{R^n}$ must contain at least $n$ vectors ($k \geq n$).
\end{theorem}


\end{frame}


\subsection{Basis and dimension in $\mathbf{R^n}$}

\begin{frame}
\frametitle{Basis and dimension in $\mathbf{R^n}$} \vskip 0.1in

It is clear from Theorem 11.6 that we can find a ``more efficient'' spanning set:

\begin{definition}
Let $\mathbf{v_1}, \ldots , \mathbf{v_k}$ be a fixed set of $k$ vectors in $\mathbf{R^n}$. Let $V$ be the set $\mathcal{L}[\mathbf{v_1}, \ldots , \mathbf{v_k}]$ spanned by $\mathbf{v_1}, \ldots , \mathbf{v_k}$. Then, $\mathbf{v_1}, \ldots , \mathbf{v_k}$ forms a \textbf{basis} of $V$ if:
\begin{itemize}
\item[(a)] $\mathbf{v_1}, \ldots , \mathbf{v_k}$ span $V$, and
\item[(b)] $\mathbf{v_1}, \ldots , \mathbf{v_k}$ are linearly independent.
\end{itemize} 
\end{definition}


\begin{example}[11.8]
 The unit vectors $$e_1 = \begin{pmatrix} 1 \\ 0 \\ \vdots \\ 0 \end{pmatrix} , \ldots , e_n = \begin{pmatrix} 0 \\ 0 \\ \vdots \\ 1 \end{pmatrix} $$ form a basis of $\mathbf{R^n}$. Since this is such a natural basis, it is called the \textbf{canonical basis}.
\end{example}

\end{frame}



\begin{frame}
\frametitle{Basis and dimension in $\mathbf{R^n}$} \vskip 0.1in

%It is clear from Theorem 11.6 that we can find a ``more efficient'' spanning set:
\begin{theorem}[11.7]
Every basis of $\mathbf{R^n}$ contains $n$ vectors.
\end{theorem}

\begin{theorem}[11.8]
Let $\mathbf{v_1}, \ldots , \mathbf{v_n}$ be a collection of $n$ vectors in $\mathbf{R^n}$. Form the $n \times n$ matrix $A$ whose columns are these $\mathbf{v}_{j}$'s: $A = \begin{pmatrix} \mathbf{v_1} &  \mathbf{v_2} & \cdots & \mathbf{v_n} \end{pmatrix}$. 

Then, the following statements are equivalent:
\begin{itemize}
\item[(a)] $\mathbf{v_1}, \ldots , \mathbf{v_n}$ are linearly independent,
\item[(b)] $\mathbf{v_1}, \ldots , \mathbf{v_n}$ span $\mathbf{R^n}$,
\item[(c)]  $\mathbf{v_1}, \ldots , \mathbf{v_n}$ form a basis of $\mathbf{R^n}$, and
\item[(d)] the determinant of $A_{n \times n}$ is nonzero
\end{itemize} 
\end{theorem}

The \textbf{dimension} of a vector space $V$ is the \# of vectors in any basis of V.

\begin{itemize}
\item[] Since every basis of $\mathbf{R^n}$ contains exactly $n$ vectors, there are $n$ independent directions in $\mathbf{R^n}$, and $\mathbf{R^n}$ is $n$-dimensional.
\end{itemize} 

\end{frame}


\section{Eigenvalues and Eigenvectors (Chapter 23)}

\subsection{Definitions}
% Ch23.1
\begin{frame}
\frametitle{Chapter 23. Eigenvalues and Eigenvectors} \vskip 0.1in

Eigenvalues (``characteristic values'') and eigenvectors of a square matrix summarize the essential properties of linear and nonlinear systems of equations:
\begin{itemize}
\item They are the components of explicit solutions to \textit{linear} dynamic models.
\item The signs of eigenvalues determine the stability of equilibria in \textit{nonlinear} dynamic models, and the definiteness of a symmetric matrix.
\item Important for economic optimization problems.
\end{itemize}

\end{frame}




%\subsection{Definitions}
% Ch23.1
\begin{frame}
\frametitle{Eigenvalues and Eigenvectors} \vskip 0.1in

\begin{example}[See 23.1, 23.2, 23.3, \& 23.4., pp. 580-1]
\end{example}

\begin{definition}
Let $A$ be a square matrix. An \textbf{eigenvalue} of A is a number $r$ which when subtracted from each of the  diagonal entries of $A$ converts $A$ into a singular matrix. Therefore, $r$ is an eigenvalue of $A$ if and only if $ A - rI $ is a singular matrix.
\end{definition}


\begin{theorem}[23.1]
The diagonal entries of a diagonal matrix $D$ are eigenvalues of $D$.
\end{theorem}

\begin{theorem}[23.2]
A square matrix $A$ is singular if and only if 0 is an eigenvalue of $A$.
\end{theorem}

A matrix is singular if and only if its determinant is zero. That is, $ A - rI $ is a singular matrix, if and only if $\label{deter} \det (A - r I) = 0 ~.$

\end{frame}


\subsection{Finding eigenvalues and eigenvectors}
% Ch23.1
\begin{frame}
\frametitle{Finding eigenvalues and eigenvectors}
For an $n \times n$ matrix $A$, $\det (A - r I)$ is an $n$th order polynomial in the variable $r$, called the \textbf{characteristic polynomial}.

\begin{example}
For a general $2 \times 2$ matrix, the characteristic polynomial is 
\begin{eqnarray} \nonumber 
\det(A - r I) &=& \det \begin{pmatrix} a_{11} - r & a_{12} \\ a_{21} & a_{22} - r\end{pmatrix} \\ \label{characpoly}
&=& r^2 - (a_{11} + a_{22})r + (a_{11}a_{22} - a_{12}a_{21}) ~,
\end{eqnarray}
a second-order polynomial. 
\end{example}
Therefore,
\begin{itemize}  
\item The eigenvalues $r$ of $A$ are the roots of the characteristic polynomial;
\item a $2 \times 2$ matrix has at most two eigenvalues (we can use the quadratic formula);  a $n \times n$ matrix has at most $n$ eigenvalues.
\end{itemize}

\end{frame}


\begin{frame}
\frametitle{Finding eigenvalues and eigenvectors}
When $r$ is an eigenvalue of $A$, a \textit{nonzero} vector $\mathbf{v}$ such that \begin{equation}\label{eigen} 
(A - r I)\mathbf{v} = \mathbf{0} 
\end{equation} 
is called an \textbf{eigenvector} of $A$ \textit{corresponding} to eigenvalue $r$. Multiplying out \eqref{eigen} yields
\begin{eqnarray}\nonumber
A\mathbf{v} - r I\mathbf{v} &=& \mathbf{0} \\  \nonumber
A\mathbf{v} - r \mathbf{v} &=& \mathbf{0} \\  \nonumber
A\mathbf{v} &=& r \mathbf{v} 
\end{eqnarray}
If $r$ is an eigenvalue and $\mathbf{v}$ is a corresponding eigenvector, then $A\mathbf{v} = r \mathbf{v}$.

\end{frame}


\begin{frame}
\frametitle{Finding eigenvalues and eigenvectors}

We can summarize the above as follows:
\begin{theorem}[23.3]\label{th23.3}
Let $A$ be an $n \times n$ matrix and let $r$ be a scalar. Then, the following statements are equivalent:
\begin{itemize}
\item[(a)] Subtracting $r$ from each diagonal entry of $A$ transforms $A$ into a singular matrix.
\item[(b)] $A - r I$ is a singular matrix.
\item[(c)] {\color{red}$\det(A - r I) = 0$}
\item[(d)] {\color{red}$(A - r I)\mathbf{v} = \mathbf{0}$} for some nonzero vector $\mathbf{v}$.
\item[(e)] $A\mathbf{v}  = r\mathbf{v}$ for some nonzero vector $\mathbf{v}$.
\end{itemize}
\end{theorem}
\end{frame}


\begin{frame}
\frametitle{Finding eigenvalues and eigenvectors}

\begin{example}
Examples 23.5 and 23.6 (pp. 583-4)
\end{example}

In general, one chooses the ``simplest'' of the nonzero candidates.
\vskip0.1in

The \textbf{eigenspace} of $A$ is the span of the set of all eigenvectors, including $\mathbf{v} = 0$. i.e., the set of all solutions to \eqref{eigen}. \vskip0.1in

In some problems, one will need to use Guassian elimination to solve the linear system $(A - r I)\mathbf{v} = \mathbf{0}$ for an eigenvector $\mathbf{v}$.

\end{frame}



\subsection{Solving linear difference equations}
% Chapter 23.2
\begin{frame}
\frametitle{Solving linear difference equations}
\framesubtitle{One-dimensional equations}
To solve a one-dimensional equation $y_{n+1} = {a}y_{n}$,

\begin{eqnarray} \nonumber
y_1 &=& {a}y_0 \\ 
\nonumber
y_2 &=& {a}y_1 = a({a}y_0) = {a^2}y_0 \\ 
\nonumber
y_3 &=& {a}y_2 = a({a^2}y_0) = {a^3}y_0 \\ \nonumber
\vdots &=& \vdots \\ \label{onedim}
y_n &=& {a^n}y_0 
\end{eqnarray} \vskip0.1in

In what settings could the solution \eqref{onedim} arise?
%% bank account; population growth

\end{frame}



\begin{frame}
\frametitle{Solving linear difference equations}
\framesubtitle{Two-dimensional systems}
Now, consider a system of two linear difference equations:
\begin{eqnarray}\label{twodim1}
x_{n+1} &=& {a}x_{n} + {b}y_{n} \\  \label{twodim2}
y_{n+1} &=& {c}x_{n} + {d}y_{n}
\end{eqnarray}
In matrix form, the system of difference equations becomes:
\begin{equation}\label{twosys}
\mathbf{z}_{n+1} = \begin{pmatrix} x_{n+1} \\ y_{n+1} \end{pmatrix} = \begin{pmatrix}
a & b \\ c & d 
\end{pmatrix} \begin{pmatrix}
x_n \\ y_n 
\end{pmatrix} \equiv A\mathbf{z}_n
\end{equation}
\begin{itemize}
\item If $b=c=0$, \eqref{twodim1} and \eqref{twodim2} are uncoupled, and can be solved as two seperate one-dimensional problems \eqref{onedim}.
\item If $b \neq 0$ or $c \neq 0$, we need to transform the coefficient matrix $A$ into a diagonal matrix so that \eqref{twodim1} and \eqref{twodim2} become uncoupled and therefore more easily solved; and
\item To transform matrix $A$ into a diagonal matrix $D$ we use the technique of a change-of-coordinates (or change-of-bases): $P^{-1}AP = D$
\end{itemize}
%% ($P^{-1}AP = D$), we use the definition of eigenvalues and eigenvectors;

\end{frame}



\begin{frame}
\frametitle{\large Visual example: conic section}
%\framesubtitle{Visual example: conic section}

\begin{figure}
\includegraphics[width=0.5\textwidth]{conicsection.jpg}
\caption{A conic section in a coordinate system adapted for it.}
\end{figure} \vskip-0.1in
\begin{eqnarray}
& Ax^2 + Bxy + Cy^2 - D = 0 & ,  B \neq 0 \nonumber \\ % A,B,C,D \in R,{~~\text{and}~}
\therefore {\text{find}}& x = \alpha{x^\prime} + \beta {y^\prime} ~,~~ y = \gamma{x^\prime} + \delta{y^\prime} & \nonumber \\
& {\text{to remove $xy$-term}} & \nonumber
\end{eqnarray}

\end{frame}



\begin{frame}
\frametitle{Solving linear difference equations}
\framesubtitle{Two-dimensional systems}
Consider the abstract system of difference equations:
$$ \mathbf{z}_{n+1} =   A\mathbf{z}_n~. $$
We want to choose a transformation $P$ and $P^{-1}$ such that:$^{*}$ 
$$\mathbf{z} = P\mathbf{Z}, ~{\text{or}}~~ \mathbf{Z} = P^{-1}\mathbf{z}~.$$
\begin{eqnarray}
\mathbf{Z}_{n+1} &=& P^{-1}\mathbf{z}_{n+1} \nonumber \\
 &=& P^{-1}(A\mathbf{z}_{n}) \nonumber \\
 &=& (P^{-1}A)\mathbf{z}_{n} \nonumber \\
 &=& (P^{-1}A)(P\mathbf{Z}_{n}) \nonumber \\
 &=& {\color{red}{(P^{-1}{A}P)}}\mathbf{Z}_{n} \nonumber \\
\mathbf{Z}_{n+1} &=& {\color{red}{D}}\mathbf{Z}_{n}  
\end{eqnarray}

\vfill
{\footnotesize{*Trace these steps in the Leslie population model example.}}
\end{frame}



%\subsection{Diagonalization}
% Ch23.1, theorem 23.4
\begin{frame}
\frametitle{Solving linear difference equations}
\framesubtitle{Diagonalization} 
It turns out \ldots since $P$ is an invertible matrix then $AP=PD$. In the two-dimensional system we can write this as: $$A[\mathbf{v}_1 ~~ \mathbf{v}_2] = [\mathbf{v}_1 ~~ \mathbf{v}_2]\begin{pmatrix}
r_1 & 0 \\ 0 & r_2 \end{pmatrix} ~,$$ where $\mathbf{v}_1$ and $\mathbf{v}_2$ are the two column vectors of the $2 \times 2$ matrix $P$. Therefore:
\begin{eqnarray}
\begin{bmatrix} A\mathbf{v}_1 & A\mathbf{v}_2
\end{bmatrix} &=& \begin{bmatrix}
r_1\mathbf{v}_1 & r_2\mathbf{v}_2
\end{bmatrix} \nonumber \\
A\mathbf{v}_1 = r_1\mathbf{v}_1~, & {\text{and}} &~~  A\mathbf{v}_2 = r_2\mathbf{v}_2 \nonumber
\end{eqnarray}
\ldots $r_1$ and $r_2$ must be eigenvalues of $A$, and $\mathbf{v}_1$ and $\mathbf{v}_2$ are the corresponding eigenvectors!! (see Theorem \href{th23.3}{23.3}) \\ \vskip0.1in


\end{frame}



\begin{frame}
\frametitle{Solving linear difference equations}
\framesubtitle{Diagonalization} 

This holds for $k$-Dimensional systems \ldots which gives us % Theorem \href{th23.4}{23.4} \ldots

\begin{theorem}[23.4]\label{th23.4}
Let $A$ be a $k \times k$ matrix. Let $r_1, r_2, \ldots, r_k$ be eigenvalues of $A$, and $\mathbf{v}_1, \mathbf{v}_2, \ldots, \mathbf{v}_k$ the corresponding eigenvectors. Form the matrix $$P = [ \mathbf{v}_1 \mathbf{v}_2 \cdots \mathbf{v}_k ] $$ whose columns are these $k$ eigenvectors. If $P$ is invertible, then 
\begin{equation}\label{diagonalisation}
P^{-1}A{P} = \begin{pmatrix}
r_1 & 0 & \cdots & 0 \\
0 & r_2 & \cdots & 0 \\
\vdots & \vdots & \ddots & \vdots \\
0 & 0 & \cdots & r_k
\end{pmatrix}
\end{equation}
Conversely, if $P^{-1}A{P}$ is a diagonal matrix $D$, the columns of $P$ must be eigenvectors of $A$ and the diagonal entries of $D$ must be eigenvalues of $A$.
\end{theorem}


\end{frame}


\begin{frame}
\frametitle{Solving linear difference equations}
\framesubtitle{A general solution}
\begin{theorem}[23.6]
Let $A$ be a $k \times k$ matrix with $k$ distinct real eigenvalues $r_1, \ldots, r_k$ and corresponding eigenvectors $\mathbf{v}_1, \ldots, \mathbf{v}_k$. The general solution of the system of difference equations $\mathbf{z}_{n+1} =   A\mathbf{z}_n$ is 
\begin{equation}\label{gensol}
\mathbf{z}_{n} =   c_{1}r_1^n\mathbf{v}_1 + c_{2}r_2^n\mathbf{v}_2 + \cdots + c_{k}r_k^n\mathbf{v}_k ~. 
\end{equation}
\end{theorem} 

\vfill 

Note: We need to know the initial vector $\mathbf{z}_0$ to solve the numerical formula $\mathbf{z}_n$ (recall \eqref{onedim}). To see this, set $n=0$ in \eqref{gensol}:
\begin{eqnarray}
\mathbf{z}_0 &=& c_{1}\mathbf{v}_1 + \cdots + c_{k}\mathbf{v}_k \nonumber \\
&=& [\mathbf{v}_1 \cdots \mathbf{v}_k] \begin{pmatrix}
c_1 \\ \vdots \\ c_k 
\end{pmatrix} = P \begin{pmatrix}
c_1 \\ \vdots \\ c_k 
\end{pmatrix}~.
\end{eqnarray} 
Therefore, for \textit{any} specific initial vector $\mathbf{z}_0$, the proper choice of $c_1, \ldots, c_k$ gives the solution for \eqref{gensol}. %[Note: this follows from the uncoupling/diagonalization of the system and easily seen in the Leslie model]
\end{frame}


\begin{frame}
\frametitle{\large An alternative approach: the powers of a matrix}
\begin{theorem}[23.7]\footnotesize
Let $A$ be a $k \times k$ matrix. Suppose that there is a nonsingular matrix $P$ such that 
\begin{equation}\nonumber
P^{-1}A{P} = \begin{pmatrix}
r_1 & 0 & \cdots & 0 \\
0 & r_2 & \cdots & 0 \\
\vdots & \vdots & \ddots & \vdots \\
0 & 0 & \cdots & r_k
\end{pmatrix} ~, ~~~~ {\text{[\eqref{diagonalisation}]}}
\end{equation}
a diagonal matrix. Then, 
\begin{equation}\nonumber
A^{n} = {P}\begin{pmatrix}
r_1^n & 0 & \cdots & 0 \\
0 & r_2^n & \cdots & 0 \\
\vdots & \vdots & \ddots & \vdots \\
0 & 0 & \cdots & r_k^n
\end{pmatrix}P^{-1} ~.
\end{equation}
The solution of the corresponding system of difference equations $\mathbf{z}_{n+1} = A\mathbf{z}_n$ with initial vector $\mathbf{z}_0$ is 
\begin{equation}\nonumber
\mathbf{z}_{n} = {P}\begin{pmatrix}
r_1^n & 0 & \cdots & 0 \\
0 & r_2^n & \cdots & 0 \\
\vdots & \vdots & \ddots & \vdots \\
0 & 0 & \cdots & r_k^n
\end{pmatrix}P^{-1}\mathbf{z}_{0} ~.
\end{equation}
\end{theorem} 

\end{frame}


% Chapter 23.2
\begin{frame}
\frametitle{Stability of equilibria} \vskip 0.1in

It follows:
\begin{itemize}
\item For $\mathbf{z}_0 = \mathbf{0}$, we have $\mathbf{z}_n  = \mathbf{0}$. Such a solution is called a \textbf{steady state}, \textbf{equilibrium}, or \textbf{stationary solution}.
\item The solution is \textbf{asymptotically stable} if every solution of $\mathbf{z}_{n+1} = A\mathbf{z}_n$ tends to the steady state $\mathbf{z} = \mathbf{0}$ as $n$ tends to infinity.
\item For every solution to have a steady state, the absolute value of all the eigenvalues of $A$ must be less than 1 $(| r_i | < 1)$ such that $r^{n}_i \to 0$ as $n \to \infty$.
\end{itemize}
\begin{theorem}[23.8]
If the $k \times k$ matrix $A$ has $k$ distinct real eigenvalues, then every solution of the general system of linear difference equations $\mathbf{z}_{n+1} = A\mathbf{z}_n$ tends to $\mathbf{0}$ if and only if all the eigenvalues of $A$ have absolute value less than 1.
\end{theorem}



\end{frame}




\subsection{Markov processes}
% Chapter 23.6
\begin{frame}
\frametitle{Markov processes} \vskip0.1in
\begin{definition}
A \textbf{stochastic process} is a rule which gives the probability that the system (or the individual in this system) will be in state $i$ at time $n + 1$, given the probabilities of its being in the various states in previous periods. \vskip0.1in

When the probability that the system is in any state $i$ at time $n + 1$ depends only on what state the system was in at time $n$, the stochastic process
is called a \textbf{Markov process}. \vskip0.1in

That is, only the \textit{immediate} past matters.
\end{definition} \vskip0.1in
The key elements of a Markov process are:
\begin{itemize}
\item[(1)] the probability $x^{i}(n)$ that state $i$ occurs at time period $n$, or alternatively, the fraction of the population under study that is in state $i$ at time period $n$; and
\item[(2)] the transition probabilities $m_{ij}$, where $m_{{\color{red}i}{\color{Green}j}}$ is the probability that the process will be in state $\color{red}i$ at time $n+1$ if it is in state $\color{Green}j$ at time $n$.
\end{itemize}

\end{frame}



\begin{frame}
\frametitle{Markov processes} \vskip0.1in

\begin{eqnarray}\label{markov}
\begin{pmatrix}
x^1(n+1) \\ \vdots \\ x^k(n+1)
\end{pmatrix} &=& \underset{\text{\textbf{Markov} (transition) \textbf{matrix}}}{\underbrace{\begin{pmatrix}
m_{{\color{red}1}1} & \cdots & m_{1k} \\
\vdots & \ddots & \vdots \\
m_{{\color{red}k}1} & \cdots & m_{kk} 
\end{pmatrix}}} \begin{pmatrix}
x^1(n) \\ \vdots \\ x^k(n)
\end{pmatrix} ~; \\ \nonumber  \\ \nonumber 
{\text{that is,~~~~~\hfill}}\mathbf{x}(n+1) &=& M\mathbf{x}(n)~,
\end{eqnarray}
where $M$ is any nonnegative matrix whose column sums $\sum_i m_{ij}$ all equal 1.\vskip0.1in

For example, each element $m_{ij}$ in first column of $M$ gives the (conditional) probability that the system will be in state ${\color{red}i=1,\ldots,k}$ next period, given that it was in state $j=1$ today. We therefore have $\sum_{i=1}^{k} m_{i1} = 1$.

\begin{example}[23.20, p. 617]
$$\begin{pmatrix}
x^{em}(n+1) \\ x^{un}(n+1)
\end{pmatrix} = \begin{pmatrix}
0.9 & 0.4 \\ 0.1 & 0.6
\end{pmatrix} \begin{pmatrix}
x^{em}(n) \\ x^{un}(n)
\end{pmatrix} $$
\end{example}

%System of equations \eqref{markov}, in which ... $M$ is a \textbf{Markov matrix} (transition matrix or stochastic matrix)) such that
\end{frame}


%p. 581 see example 23.4
\begin{frame}
\frametitle{Markov processes} \vskip0.1in
Some general principles that example 23.20 illustrates
\begin{theorem}[23.15]
Let $M$ be a positive Markov matrix. Then,
\begin{itemize}
\item[(a)] $r_1 = 1$ is an eigenvalue of every $M$;
\item[(b)] every other eigenvalue $r$ of $M$ satisfies $|r|<1$
\item[(c)] eigenvalue $r = 1$ has an eigenvector $\mathbf{w}_1$ with strictly positive components; and
\item[(d)] if we write $\mathbf{v}_1$ for $\mathbf{w}_1$ divided by its components, then $\mathbf{v}_1$ is a probability vector and each solution $\mathbf{x}(n)$ of $\mathbf{x}(n+1) = M\mathbf{x}(n)$ tends to $\mathbf{v}_1$ as $n \to \infty$.
\end{itemize}
\end{theorem}
\end{frame}

%exercises 23.30 and 23.33


\subsection{Symmetric matrices}
% Chapter 23.7
\begin{frame}
\frametitle{Symmetric matrices} \vskip 0.1in

Most matrices that arise in optimization and econometrics are symmetric matrices (e.g., Hessians). If $A$ is a symmetric matrix. Then,
\begin{definition}
\begin{itemize}
\item[] $A = A^T$;
\item[] $A$ is a $k \times k$ square matrix;
\item[] its ``counter-diagonals'' have the same entries: $a_{ij} = a_{ji}$ ;
\item[] it is orthogonally diagonalizable, meaning that
we can find an orthogonal matrix $P$ ($P^{-1}=P^T$) which diagonalizes $A$: $D = P^{-1}AP$.
\end{itemize}
\end{definition}\vfill

See also Section 23.8: Definiteness
\end{frame}







%\begin{frame}
%\begin{center}
%\Huge{THE END}
%\end{center}
%\end{frame}


\end{document}
